\section{Physical validation: spin scaling, pole census, sum rule and size}
\label{sm:physics}

\paragraph{The charge scale, in full.} $\Delta(L)$ decreases monotonically and stays above the Lieb--Wu
thermodynamic value $4.67952\,t$ at every size, the residual falling from $0.679\,t$ at $L=6$ to $0.289\,t$ at
$L=12$---that is, to $6.2\%$ of the thermodynamic value (Table~\ref{tab:gaps}).
$\Delta(L=12)=4.96876\,t$ reproduces the $4.97\,t$ this work has quoted throughout, and is the same
number that appears as the Hubbard-band separation of Fig.~\ref{fig:lattice} and as the low-frequency
scale of Fig.~\ref{fig:sqw}. We fit no extrapolation exponent: a straight line in $1/L$ through the
four points extrapolates to $4.561\pm0.044\,t$, which undershoots Lieb--Wu by $0.12\,t$ and is
therefore evidence that the functional form is not $1/L$, not evidence about the limit. The statement
supported is that the charge scale converges to a finite thermodynamic value and is already within
$6.2\%$ of it at the largest size reached here.

\paragraph{Spin.} The lowest pole of $S^{zz}(q=\pi,\omega)$ gives $L\,\Delta_{s}=2.091$, $1.917$,
$1.996$, $1.996$: a $1/L$ closing with a coefficient of about $2$ (mean $2.000$, range $1.917$--$2.091$). Two fits make that precise and
bound it honestly. Fitting $\Delta_{s}=c_0+c_1/L$ over these four sizes gives an intercept
$c_0=-0.020\pm0.020$---consistent with zero at one standard error---against
$c_0=4.561\pm0.044$ for the charge gap: one intercept is consistent with zero and the other is not,
and that separation, not either intercept on its own, is the result. It has to be read with the control
that limits it, which we state here rather than leave to the figure: the pure Heisenberg ring, whose
spin gap is known to vanish in the thermodynamic limit, gives a linear-in-$1/L$ intercept of
$0.0142\pm0.0014\,t$---above zero, not consistent with it---so at these sizes a vanishing intercept is
not by itself evidence that a gap closes. The statement that survives without a fitted model is the
flatness of $L\,\Delta_s$ below. Fitting instead $L\,\Delta_{s}=d_0+d_1L$
gives $d_0=2.09\pm0.17$ with $d_1=-0.010\pm0.018$, a slope consistent with zero: over $L=6$--$12$ the
product $L\,\Delta_s$ is flat to the resolution of four points. What is \emph{not} supported is a
third significant figure on the coefficient. The four values span $8\%$ and are not monotone, and we
quote $L\,\Delta_{s}\approx2$ with the spread attached, not $2.00$; the intercept $d_0=2.09$ is the fit's value at $L=0$, not the level over $L=6$--$12$.

The independent control is the effective spin model. The pure Heisenberg ring at $J=4t^{2}/U=0.5\,t$
---no charge degree of freedom, no double occupancy, a different Hamiltonian on a different Hilbert
space, solved by dense diagonalization on a different code path---gives
$\Delta_{\rm ST}(N=12)=0.35585\,J=0.17792\,t$ and reproduces the same $1/L$ law across the same four
sizes, with $L\,\Delta_{\rm ST}=2.054$, $2.091$, $2.116$, $2.135$. Two things have to be said about
that agreement rather than left for a referee. First, the two numbers at $L=12$ differ by $6.5\%$ of
the Heisenberg value, equivalently $7.0\%$ of the Hubbard one; we name the denominator because the
two differ in the third figure. Second, the residual is not monotone and changes sign---the Hubbard
gap is $1.8\%$ \emph{above} the Heisenberg one at $L=6$ and $8.3\%$, $5.7\%$ and $6.5\%$
\emph{below} it at $L=8,10,12$---so it cannot be attributed to a smooth $O(t^{4}/U^{3})$ correction
to the effective exchange, which would not change sign. A correction of that order is the expected
\emph{scale} of the difference at $U/t=8$, and it is the right scale; its size dependence is not
resolved here and we claim nothing about it. The control confirms the $1/L$ scaling and the magnitude
of the coefficient to within $10\%$; it does not pin the coefficient, and its own $L\,\Delta_{\rm ST}$
drifts upward by $0.0134\pm0.0014$ per site over the same range.

What establishes the physics instead is measurement rather than illustration: the four-size scaling of
Table~\ref{tab:gaps}, plotted in Fig.~\ref{fig:gapscaling}, in which the spin scale is a converged
pole position at every size and closes as
$\approx2/L$ with an independent effective-Heisenberg control, and the pole census of
Table~\ref{tab:poles}, which states what is actually inside the window at each size instead of naming
it.

\subsection{A pole census}

At $L=12$ and the broadenings used here, none of the three momentum-resolved channels is a dense
spectrum inside the plotted window; each momentum carries a small number of discrete poles that the
broadening merges into a band. Counting poles that carry more than $1\%$ of the channel
maximum---Lanczos with full reorthogonalization at $n_\ell=120$, checked for depth stability at
$n_\ell=60$, $120$ and $240$ at $q=\pi$ in every channel, and validated pole by pole against dense
diagonalization at $L=6$ and $L=8$~\cite{haydock1972,golub2010}---gives Table~\ref{tab:poles}. The
census is the deposited \texttt{pole\_structure.json} and is taken on the momentum grid the figures of
this section plot, seven values of $k$ and six of $q$; a count over a coarser grid is a count of
different momenta and is not comparable entry by entry. The
census covers the three channels the figures of this section plot; the $k$-integrated $A(\omega)$,
which is a sum over the same branch, and the removal branch are not counted, and no statement is made
about them here.

\begin{table*}[tp]\centering\small
\caption{\textbf{Pole census, across three sizes.} Poles carrying more than $1\%$ of the channel
maximum, counted per momentum; each entry is the range over the momenta plotted ($7$ values of $k$ for
$A(k,\omega)$, $6$ values of $q$ for the collective channels). The $L=6$ and $L=8$ columns are dense
diagonalizations, where the count needs no recursion, on the subset of momenta available at those
sizes; the $L=12$ column is the reorthogonalized recursion at $n_\ell=120$, which reproduces the dense
counts pole by pole at the two smaller sizes. ``Leading weight'' is the fraction carried by the single
largest pole at $L=12$. The addition branch of $A(k,\omega)$ and the charge channel are at
$\eta=0.18\,t$, the spin channel at $\eta=0.10\,t$. The last column is the smallest separation between
retained poles in units of the full width at half maximum, which is why the figures show bands: at the
tightest momenta the lines are genuinely merged, and the census says how many of them there
are.}\label{tab:poles}
\setlength{\tabcolsep}{4pt}
\begin{tabular}{lcccccc}
\toprule
 & \multicolumn{3}{c}{poles above $1\%$} & & leading & $\delta\omega_{\min}/$\\
\cmidrule(lr){2-4}
channel & $L{=}6$ & $L{=}8$ & $L{=}12$ & $\dim$ ratio & weight & FWHM\\
\midrule
$A(k,\omega)$, addition  & $2$--$3$ & $2$--$4$ & $3$--$7$ & $2440$ & $0.32$--$0.69$ & $0.12$--$2.04$\\
$S(q,\omega)$, charge    & $3$--$4$ & $3$--$6$ & $6$--$9$ & $2440$ & $0.38$--$0.68$ & $0.24$--$2.20$\\
$S^{zz}(q,\omega)$, spin & $1$--$2$ & $1$--$3$ & $1$--$3$ & $2134$ & $0.80$--$0.99$ & $1.08$--$2.37$\\
\bottomrule
\end{tabular}
\end{table*}

The count does grow with size, and the size of the growth is the point. Between $L=6$ and $L=12$ the
$(N{+}1)$ sector grows by a factor $2440$, while the largest per-momentum pole count grows from $3$ to
$7$ in the addition branch, from $4$ to $9$ in the charge channel and from $2$ to $3$ in the spin
channel. A line count that roughly doubles while the Hilbert space grows by three
orders of magnitude is a discrete spectrum whose resolvable lines are being broadened together, not a
dense one forming. The spin channel is the extreme case: at every momentum a single pole carries
between $80\%$ and $99\%$ of the weight, at separations of at least $1.08$ full widths, and at $q=\pi$
that single pole \emph{is} the $0.16635\,t$ scale of Table~\ref{tab:gaps}. What Fig.~\ref{fig:spin}
shows at $L=12$ is one dominant antiferromagnetic mode plus at most two satellites.

This has a consequence for how the reference itself must be described: $A(k,\omega)$ is not computed exactly, as the Lehmann sum would be. With $n_\ell=150$ and no reorthogonalization in a sector of dimension $731\,808$, what
is computed is a $150$-node Gauss quadrature of the spectral measure~\cite{haydock1972,golub2010}. It
is accurate enough here---and it is accurate \emph{precisely because} the measure it integrates is a
handful of poles; a dense spectrum at this dimension would not be reproduced by $150$ nodes. It is not,
however, converged: the same depth moves the reference by $0.93\%$ at $\eta=0.15\,t$ and by $3.0\%$ at
the $\eta=0.10\,t$ of Fig.~\ref{fig:spin} when it is taken from $n_\ell=100$ to $n_\ell=200$
(Sec.~\ref{sec:prices}). The reference is therefore labelled for what it is: a finite-depth continued
fraction, checked against dense diagonalization at $L\le8$ and quoted with its depth. The pole
positions of Table~\ref{tab:gaps} are a separate and stronger object---they are computed at full
reorthogonalization and are stable across $n_\ell=60$, $120$, $240$---which is why the physical claim
of Sec.~\ref{sec:gaps} rests on them and not on a lineshape. The census counts poles, and the count
does not depend on the broadening either; only the picture does.

\subsection{The sum rule, as measured}

The analytic identity $\int A^{+}\,d\omega = 1-\langle\hat n_p\rangle = 0.5$ is exact and
grid-independent, but it is an identity of the \emph{ground state}: evaluating it costs one
ground-state expectation value, $\langle\Psi_0|\hat c_p\hat c^{\dagger}_p|\Psi_0\rangle$, and never
touches the reconstruction. It is therefore not an a posteriori check of $A_{\mathcal S}$, and the caption of Fig.~\ref{fig:method} does not present it as one.

Integrating the plotted curves instead, over the $600$-point window at $\eta=0.15\,t$:
\begin{widetext}
\begin{equation}
\int_{\rm window} A_{\rm exact}\,d\omega = 0.486588,
\qquad
\int_{\rm window} A_{\mathcal S}\,d\omega = 0.486471 .
\end{equation}
\end{widetext}
Both fall $2.7\%$ below the analytic $0.5$ because Lorentzians of width $\eta$ leak outside any finite
window: the analytic in-window mass is $0.486589$, the tail outside it is $0.013411$, and the residue
of the trapezoidal quadrature itself is $3\times10^{-6}$ relative---so the deficit is entirely the
window and not the reconstruction. The a posteriori statement the reconstruction does support is the
comparison of the two integrated curves, which differ by $1.16\times10^{-4}$ absolute, that is
$2.4\times10^{-4}$ relative. That is the number we report, and it is consistent with the ${\sim}2\%$ window effect that Fig.~\ref{fig:lattice} states.

\subsection{Size: the convention, and what it does not license}

The system size is not hidden and is not the weak point. The exact-diagonalization convention for this
observable is $L=12$--$16$ at $\eta\approx0.1$--$0.2\,t$: Xu, Lu, Tohyama and Shao~\cite{xu2024}
compute this spectral function by Lanczos at $L=12$ and $\eta=0.2\,t$ at half filling, and at
$L=16$ at quarter filling. Their model is the extended Peierls--Hubbard chain, and their $L=16$ run is
a dimerized lattice away from half filling; it is their half-filled $L=12$, and the plain-Hubbard
$L=12,8,6$ of their Appendix~A, that sit alongside ours. Our $L=12$ at $\eta=0.10$--$0.18\,t$ is that convention, not an outlier
within it.

What the convention does not license is the vocabulary of the thermodynamic limit, and the criterion
that says so is in the dynamical-DMRG literature this paper already cites. Jeckelmann's rule is
explicit~\cite{jeckelmann2002}: \emph{``$\eta_0(N)$ must be larger than the distance $\delta\omega$
between two consecutive excited states.''} At $L=12$ the largest spacing between consecutive
excited-state energies in the window is ${\approx}0.48\,t$, a factor $2.7$ \emph{above} the
$\eta=0.18\,t$ used for $A(k,\omega)$ and $4.8$ above the $\eta=0.10\,t$ used for the spin channel.
The rule is not satisfied here. Satisfying it would demand $\eta>0.48\,t$. That is only a tenth of
the Mott scale, so the Hubbard bands would survive; but it is $0.96\,J$ and $2.9$ times the $L=12$
spin gap itself, so the entire low-frequency structure of Fig.~\ref{fig:spin} would merge into one
feature. The number the rule
demands is set by the spacing, $0.48\,t$; we take it from the spacing and not from a rounded-up
interval. The correct response is not a larger $\eta$: it is to stop using words that require
the rule to hold. The dynamical-DMRG standard fixes the product $\eta L$ and lets $\eta$ fall with
the size: Benthien and Jeckelmann use $\eta=5t/L$ in one place, the $L=100$ spin calculation, and
declare $\eta L=12\,t$ for the finite-size scaling that carries the charge channel to
$L=120$~\cite{benthien2007}; Jeckelmann's own prescription is $\eta(N)=c/N$ with $c=6.4$ and
$12.8\,t$, on chains of up to $N=256$~\cite{jeckelmann2002}. This work is at $L=12$ with
$\eta L=1.2$--$2.2\,t$, a factor of two to ten below every constant in that range,
and it is the words, not the size, that we withdraw. Section~\ref{sec:scope:declared} states the same
concession in the list of what this work does not claim, so that it is not read only here.
